% SLIDE 17: Key Takeaways
\begin{frame}{Key Takeaways}
\centering

\begin{tikzpicture}[node distance=1.0cm]
  \node[pipeblock=garnet, minimum width=5cm, minimum height=1.0cm] (c1) at (0,2) {
    \textbf{1. CSTM} --- canonical camera transform
  };
  \node[pipeblock=atlantic, minimum width=5cm, minimum height=1.0cm] (c2) at (0,0.5) {
    \textbf{2. Joint optimization} --- depth + normals
  };
  \node[pipeblock=congaree, minimum width=5cm, minimum height=1.0cm] (c3) at (0,-1.0) {
    \textbf{3. RPNL} --- local depth normalization
  };
\end{tikzpicture}

\vspace{0.5cm}
\textbf{Big idea:} Transforming the \textit{data space} (CSTM) is simpler and more effective than encoding camera parameters into the network architecture.

\vspace{0.3cm}
Geometric foundation models trained on massive diverse data can rival --- and often beat --- task-specific models fine-tuned on individual benchmarks.

\vspace{0.5cm}
{\large\textcolor{garnet}{Questions?}}
\end{frame}
